\begin{frame}[t]{UniPert--G2CP：靶点表示与两类响应数据}
\bodyfont\renewcommand{\arraystretch}{1.22}
\begin{tabularx}{\textwidth}{@{}P{2.6cm}Y Y@{}}
\toprule
\textbf{资源 / 模块}&\textbf{正式版报告的规模}&\textbf{测量或输入类型}\\\midrule
UniPert：PubChem CPI & 53,963 个互作；8,533 个化合物；4,250 个蛋白靶点（2024-01 快照） & 化合物结构、靶点关系及蛋白序列；不是实验蛋白丰度\\
G2CP：LINCS & A375、A549、HT29、MCF7、PC3；978 个 landmark 转录本 & Level 5 转录组签名；CRISPR 96 h，化合物 10 $\mu$M / 24 h\\
G2CP：CPJUMP1 & A549、U2OS；160 个遗传、302 个化学扰动；838 维特征 & well-level 细胞形态特征，由图像计算获得\\
\bottomrule
\end{tabularx}
\vspace{.22cm}
\subhead{与本项目的关系}
可复用扰动编码器与迁移流程；需要另行引入质谱蛋白矩阵。\textbf{蛋白序列编码不等于蛋白丰度监督。}验证集还含 Dixit / Norman / sciPlex3，分项规模见讲稿。
\source{[22] Results、STAR Methods。LINCS 精确逐细胞系条数见独立 Table S4，本次未取得。正式版表示空间分析为 4,994 基因 / 7,860 化合物。}
\end{frame}
